%# -*- coding: utf-8-unix -*-

\chapter{RATIONAL APPROXIMATIONS TO ALGEBRAIC NUMBERS\index{numbers}}\label{ch:7}

\section{Introduction}\label{sec:ch7_1}

In 1909, a remarkable improvement on Liouville\index{Liouville, J.}'s theorem\index{Liouville's theorem} was obtained by the Norwegian mathematician Axel Thue\index{Thue, A.}.\footnote{J.M. 135 (1909), 284--305.} He proved that for any algebraic number\index{algebraic number}  $\alpha$  with degree n > 1 and for any  $\kappa > \frac{1}{2}n + 1$  there exists  $c = c(\alpha, \kappa) > 0$  such that  $|\alpha - p/q| > c/q^{\kappa}$  for all rationals p/q (q>0). His work rested on the construction of an auxiliary polynomial in two variables possessing zeros to a high order, and it can be regarded as the source of many of our modern transcendence techniques. The condition on  $\kappa$  was relaxed by Siegel\index{Siegel, C. L.}\footnote{M.Z. 10 (1921), 173--213.} in 1921 to  $\kappa > s + n/(s+1)$  for any positive integer s, thus, in particular, to  $\kappa > 2\sqrt{n}$, and it was further relaxed by Dyson\index{Dyson, F.}\footnote{Acta Math. 79 (1947), 225--40.} and Gelfond\index{Gelfond, A. O.}\footnote{See Bibliography.} independently in 1947 to  $\kappa > \sqrt{2n}$. The latter expositions continued to involve polynomials in two variables and further progress seemed to require some extension of the arguments relating to polynomials in many variables; in fact special results in this connexion had already been obtained by Schneider\index{Schneider, Th.}\footnote{J. M. 175 (1936), 182--92.} in 1936. A generalization of the desired kind was discovered by Roth\index{Roth, K. F.}\footnote{Mathematika, 2 (1955), 1--20.} in 1955; he showed indeed that the above proposition holds for any  $\kappa > 2$, a condition which, in view of the introductory remarks of \autoref{ch:1}, is essentially best possible.

Roth\index{Roth, K. F.}'s work, however, gave rise to a number of further problems. Siegel\index{Siegel, C. L.} had initiated studies on the approximation of algebraic numbers\index{algebraic number} by algebraic numbers\index{algebraic number} in a fixed field, and also by algebraic numbers\index{algebraic number} with bounded degree, and although Roth\index{Roth, K. F.}'s arguments could be readily generalized to furnish a best possible result in connexion with the first topic,\footnote{See LeVeque\index{LeVeque, W. J.} (Bibliography).} they did not seem to admit a similar extension in connexion with the second. Even less, therefore, did they appear capable of dealing with the wider question concerning the simultaneous approximation of algebraic numbers\index{algebraic number} by rationals. The whole subject was resolved by Schmidt\index{Schmidt, W. M.}\footnote{See Bibliography.} in 1970; building upon Roth\index{Roth, K. F.}'s foundations but introducing several new ideas, in particular from the Geometry of Numbers\index{Geometry of Numbers}, he proved:

\begin{mythm}{}{ch7_1}
For any algebraic numbers\index{algebraic number}  $\alpha_1, ..., \alpha_n$  with  $1, \alpha_1, ..., \alpha_n$  linearly independent over the rationals, and for any  $\epsilon > 0$, there are only finitely many positive integers q such that
\[q^{1+\epsilon}\|q\alpha_1\|\dots\|q\alpha_n\|<1.\]
\end{mythm}

Here ||x|| denotes the distance of x from the nearest integer taken positively. The theorem implies, by a classical transference principle,\footnote{See Cassels\index{Cassels, J. W. S.}' Diophantine approximation (Bibliography).} that there are only finitely many non-zero integers  $q_1, ..., q_n$  with
\[|q_1 \dots q_n|^{1+\epsilon} ||q_1 \alpha_1 + \dots + q_n \alpha_n|| < 1.\]

Further, as immediate corollaries, we see that there are only finitely many integers $p_1, \ldots, p_n, q \ (q > 0)$  satisfying
\[|\alpha_j - p_j/q| < q^{-1-(1/n)-\epsilon} \quad (1 \leqslant j \leqslant n),\]
and also only finitely many integers $p, q_1, ..., q_n$  satisfying
\[|q_1\alpha_1+\ldots+q_n\alpha_n-p|< q^{-n-\epsilon},\]
where  $q = \max |q_j|$. Furthermore we have:

\begin{mythm}{}{ch7_2}
For any algebraic number\index{algebraic number}  $\alpha$  with degree exceeding n and any  $\epsilon > 0$, there are only finitely many algebraic numbers\index{algebraic number}  $\beta$  with degree at most n such that  $|\alpha - \beta| < B^{-n-1-\epsilon}$, where B denotes the height of  $\beta$.
\end{mythm}

The theorem follows from the inequality just above with  $\alpha_j = \alpha^j$, on noting that, if P(x) is the minimal polynomial\index{minimal polynomial of algebraic number} for  $\beta$, then
\[|P(\alpha)| < BC|\alpha - \beta|\]
for some C depending only on  $\alpha$. The exponent of B is essentially best possible, as has been demonstrated by Wirsing\index{Wirsing, E.}. In fact, Wirsing\index{Wirsing, E.} obtained \autoref{thm:ch7_2} in 1965 before the work of Schmidt\index{Schmidt, W. M.}, but with the less precise exponent  $-2n - \epsilon$.\footnote{J.M. 206 (1961), 67--77.}

One of the main applications of the methods of this chapter has concerned Diophantine equations\index{Diophantine equations} of norm form in several variables, which generalize the Thue\index{Thue, A.} equation discussed in \autoref{ch:4}; indeed the work has led to a complete description of all such equations that possess only finitely many solutions.\footnote{Proc. Symposia Pure Math. (Amer. Math. Soc.), 20 (1971), 213--47.}

\begin{mythm}{}{ch7_3}
Let K be an algebraic number\index{algebraic number} field and let M be a module in K. A necessary and sufficient condition for there to exist an integer m such that the equation  $N\mu = m$  has infinitely many solutions  $\mu$  in M is that M be a full module in some subfield of K which is neither the rational nor an imaginary quadratic field.
\end{mythm}

The necessity follows at once from the fact that the subfield, if it exists, contains at least one fundamental unit, and the sufficiency is a consequence of a generalized version of \autoref{thm:ch7_1} relating to products of linear forms;\footnote{M.A. 191 (1971), 1--20.}  it is in fact a direct corollary in the case when the dimension of M is small compared with the degree of K. As examples, one sees that the equation
\[N(x_1 + x_2\sqrt{2} + x_3\sqrt{3}) = 1\]
has infinitely many solutions in integers  $x_1$, $x_2$, $x_3$  given by
\[x_1 + x_2\sqrt{2} = \pm (1 + \sqrt{2})^n, \quad \text{and by} \quad x_1 + x_3\sqrt{3} = \pm (2 + \sqrt{3})^n,\]
where n = 0, 1, 2, ...; and the equation
\[N(x_1+q^{1/p}x_2+\ldots+q^{(p-2)/p}x_{p-1})=m,\]
where p, q are primes and m is any integer, has only a finite number of solutions in integers  $x_1, ..., x_{p-1}$; for clearly the field generated by  $q^{1/p}$  over the rationals has only trivial subfields. It should be noted, however, that, in contrast to the work of \autoref{ch:4}, the arguments here are not effective and cannot lead to a determination of the totality of solutions. In fact, apart from a few special results of Skolem\index{Skolem, Th.},\footnote{See Bibliography.} the only effective theorems established to date on equations of norm form in three or more variables derive from the work on the hypergeometric function\index{hypergeometric function} referred to in \S~5 of \autoref{ch:4}.

A generalization of Roth\index{Roth, K. F.}'s theorem\index{Roth's theorem} in the p-adic domain was obtained by Ridout\index{Ridout, D.}\footnote{Mathematika, 4 (1957), 125--31; 5 (1958), 40--8; see also Mahler\index{Mahler, K.} (Bibliography).} in 1957; in particular he proved that for any algebraic number\index{algebraic number}  $\alpha$  and any  $\epsilon > 0$, there exist only finitely many integers p, q, comprised solely of powers of fixed sets of primes, such that  $|\alpha - p/q| < q^{-\epsilon}$. In this case, however, \autoref{thm:ch3_1}\footnote{P.C.P.S. 63 (1967), 693--702.} gives rather more; in fact, on taking  $\alpha_1 = \alpha$  and the remaining  $\alpha$'s as the given primes, one sees at once that  $q^{-\epsilon}$  can be replaced by  $(\log q)^{-c}$  for some c depending only on  $\alpha$  and the primes, a result moreover that is fully effective. Further theorems in the context of p-adic approximations follow from the other inequalities for  $|\Lambda|$  recorded in \autoref{ch:3}.

\section{Wronskians}\label{sec:ch7_2}

The Wronskian of polynomials  $\phi_1(x), \ldots, \phi_k(x)$  of one variable is defined as the determinant of order k with  $(j!)^{-1}\phi_i^{(j)}(x)$  in the ith row and (j+1)th column, where  $1 \leq i \leq k, \ 0 \leq j < k$, and  $\phi^{(j)}$  denotes the jth derivative of  $\phi$. Such Wronskians\index{Wronskians} occurred in the original work of Thue\index{Thue, A.}, and they sufficed for the expositions of Siegel\index{Siegel, C. L.}, Dyson\index{Dyson, F.} and Gelfond\index{Gelfond, A. O.}; the arguments of Roth\index{Roth, K. F.} and Schmidt\index{Schmidt, W. M.}, however, involved the concept of a generalized Wronskian. Suppose that  $\phi_1, \ldots, \phi_k$  are polynomials in n variables  $x_1, \ldots, x_n$  and let  $\Delta^{(j)}$  denote a differential operator of the form
\[(j_1!\ldots j_n!)^{-1}(\partial/\partial x_1)^{j_1}\ldots(\partial/\partial x_n)^{j_n},\]
where  $j_1 + \ldots + j_n = j$. Then any determinant of order k with some  $\Delta^{(j)}\phi_i$  in the ith row and (j+1)th column is called a generalized Wronskian of  $\phi_1, \ldots, \phi_k$. There are clearly only finitely many generalized Wronskians\index{Wronskians} of  $\phi_1, \ldots, \phi_k$, and when n=1 the set reduces to the original Wronskian. We shall require later the result that if  $\phi_1, \ldots, \phi_k$  are linearly independent over their field of coefficients then some generalized Wronskian does not vanish identically; proofs are given, for instance, in the tracts of Cassels\index{Cassels, J. W. S.} and Mahler\index{Mahler, K.}.

\section{The index}\label{sec:ch7_3}

The proof of \autoref{thm:ch7_1} involves polynomials P in kn variables  $x_{lm}$  ( $1 \le l \le k$,  $1 \le m \le n$ ), homogeneous in  $x_{1m}, \ldots, x_{km}$  for each m. Suppose that P has real coefficients and let  $L_m$  ( $1 \le m \le n$ ) be real linear forms in  $x_{1m}, \ldots, x_{km}$. Then the index of P with respect to  $L_1, \ldots, L_n$  and positive integers  $r_1, \ldots, r_n$  is defined as the largest value of  $\theta$  such that P can be expressed as a linear combination of the polynomials  $L_1^{j_1} \ldots L_n^{j_n}$  with
\[(j_1/r_1) + \ldots + (j_n/r_n) \geq \theta,\]
and with real polynomials in the  $x_{lm}$  as coefficients. It is easily verified that, for any polynomials P, Q as above, the index, for brevity, ind, with respect to the  $L_m$  and  $r_m$  satisfies ind  $(P+Q) \ge \min (\operatorname{ind} P, \operatorname{ind} Q)$,
\[\operatorname{ind} PQ = \operatorname{ind} P + \operatorname{ind} Q.\]

We shall require also the related concept of the index of a real polynomial  $P(x_1, ..., x_n)$  with respect to rationals  $p_m/q_m$   $(q_m > 0)$  and integers  $r_m > 0$   $(1 \le m \le n)$; this is defined as the index of the polynomial  $x_{01}^{d_1} \dots x_{0n}^{d_n} P(x_{11}/x_{21}, \dots, x_{1n}/x_{2n})$  in the 2n variables  $x_{lm}$  (l = 1, 2) with respect to the linear forms
\[L_m = q_m x_{1m} - p_m x_{2m}\]
and the  $r_m$, where  $d_m$  denotes the degree of P in  $x_m$. The index in the latter sense occurred first in the work of Roth\index{Roth, K. F.}, and the generalized concept was introduced by Schmidt\index{Schmidt, W. M.}.

In analogy with the notation of earlier chapters, we define the height ||P|| of a polynomial P as the maximum of the absolute values of its coefficients; we shall speak of the height only for polynomials with rational integer coefficients, not identically 0. The same definition will of course apply in the special case of linear forms.

Suppose now that P is a polynomial in kn variables as indicated at the beginning of the section. Let  $L_1, \dots, L_n$  be linear forms as there, with relatively prime integer coefficients, and let  $q_m = \|L_m\|$. Further let  $r_1, \dots, r_n$  be positive integers such that  $\delta r_m > r_{m+1}$   $(1 \le m < n)$, where  $\delta = (\epsilon/32)^{2^n}$  and  $0 < \epsilon < 1$. We have

\begin{mylemma}{}{ch7_1}
If  $q_m^{r_m} > q_1^{r_{m+1}} \ (1 \le m < n)$  and  $q_1^{\delta \eta} > 8^{nk^2}$, where  $0 < \eta \le k$, and if also P has height at most  $q_1^{\delta \eta r_1/k^2}$  and degree at most  $r_m$  in  $x_{1m}, \ldots, x_{km}$, then the index of P with respect to the  $L_m$  and  $r_m$  is at most  $\epsilon$.
\end{mylemma}

This is an extension, due to Schmidt\index{Schmidt, W. M.}, of the most fundamental part of Roth\index{Roth, K. F.}'s work, sometimes called Roth's lemma\index{Roth's lemma}. The result follows easily in fact from the case considered by Roth\index{Roth, K. F.}, as we now show.

We assume, as we may without loss of generality, that  $q_m = |a_{1m}|$, where
\[L_m = \sum_{l=1}^k a_{lm} x_{lm} \quad (1 \leqslant m \leqslant n).\]

We shall further assume that $(a_{1m}, a_{2m})$\footnote{(a, b) denotes the greatest common divisor of a, b.} is at most $q_m^{(k-2)/(k-1)}$; this also involves no loss of generality, since a prime $p$ can divide at most $k-2$ of the integers $(a_{1m}, a_{lm})$ with $1 < l \le k$, whence their product divides $q_m^{k-2}$. Let now $P'$ be the polynomial obtained from $P$ by successively removing, in some order, the highest power of
\[x_{lm} \quad (3 \leqslant l \leqslant k, 1 \leqslant m \leqslant n)\]
that divides $P$ and then setting the variable to 0; further let $P''$ be the polynomial obtained by setting $x_{2m} = 1$ in $P'$ for each $m$. Then clearly the index of $P$ with respect to the $L_m$ and $r_m$ is at most the index of $P''$ with respect to $-a_{2m}/a_{1m}$ and $r_m$. Also, by assumption, the denominator of $a_{2m}/a_{1m}$, when expressed in lowest terms, namely $q_m/(a_{1m}, a_{2m})$, is at least $q_m^{1/k}$. Hence we see that it suffices to prove the following modified version of \autoref{lem:ch7_1}.

For any integers $r_m$ $(1 \leq m \leq n)$ as above and any rationals
\[p_m/q_m \quad (q_m > 0)\]
in their lowest terms such that $q_m^{r_m} > q_1^{r_{m+1}}$ and $q_1^{\delta \eta} > 8^n$, where $0 < \eta \leqslant 1$, the index with respect to the $p_m/q_m$ and $r_m$ of any polynomial $P(x_1,\dots,x_n)$ with height at most $q_1^{\delta \eta r_1/k^2}$ and degree at most $r_m$ in $x_m$ is at most $\epsilon$.

Proofs of this proposition, possibly in slightly adapted form, in particular with $\eta=1$, are given in several of the texts cited in the Bibliography, and our exposition can therefore be relatively brief. The result plainly holds for $n=1$, for if $j_1$ is the exponent to which $x_1-p_1/q_1$ divides $P(x_1)$ then, by Gauss\index{Gauss, C. F.}' lemma\index{Gauss' lemma}, we have
\[P(x_1) = (q_1x_1 - p_1)^{j_1}Q(x_1),\]
where $Q$ is a polynomial with integer coefficients; thus the leading coefficient of $P$ is at least $q_1^{j_1}$, whence $j_1/r_1 < \delta \eta < \epsilon$, as required. We now assume the validity of the proposition with $n$ replaced by $n-1$ and we proceed to establish the assertion for $n \ (\geq 2)$.

We begin by writing $P$ in the form
\[\phi_0 \psi_0 + \dots + \phi_{s-1} \psi_{s-1}\]
where the $\phi$'s and $\psi$'s are polynomials in the variables $x_1, \dots, x_{n-1}$ and $x_n$ respectively with rational coefficients, and we choose one such representation for which $s \ (\leq r_n + 1)$ is minimal. Then there exist Wronskians\index{Wronskians} $U'$, $V'$ of the $\phi$'s and $\psi$'s respectively which do not vanish identically, and clearly $W = U'V'$ can be expressed as a determinant of order $s$ with $\Delta^{(j)}(i!)^{-1} (\partial/\partial x_n)^i P(x_1, \dots, x_n)$ in the $(i+1)$-th row and $(j+1)$-th column, where the $\Delta^{(j)}$ are operators as in \S~2 with $j_n = 0$. Hence $W$ is a polynomial with degree at most $s r_j$ in $x_j$ and with $\|W\| \leq (8^{rn} \|P\|)^s \leq q_1^{2\delta\eta r_1 s}$, where $r=r_1=\max r_m$; here we are using the hypothesis $q_1^{\delta\eta}>8^n$ and the observations that $\Delta^{(j)}$ acting on any monomial in $P$ introduces a factor not exceeding $2^{rn}$, that there are at most $2^{rn}$ such monomials, and that the number of terms obtained on expanding the determinant for $W$ is $s!\leqslant 2^{rs}$. Now, again by Gauss\index{Gauss, C. F.}' lemma\index{Gauss' lemma}, we have $W=UV$, where $U,V$ are polynomials with integer coefficients in the variables $x_1,\dots,x_{n-1}$ and $x_n$ respectively, given by some rational multiples of $U',V'$; and clearly the bound for $\|W\|$ obtains also for $\|U\|$ and $\|V\|$. Thus, by our inductive hypothesis, it follows, on taking $2\delta$ in place of $delta$, that the index of $U$ with respect to the $p_m/q_m$ and $r_m$ is at most $2^{-5+1/2^{n-1}} \epsilon^2 s$. Further, by the case $n=1$ of the proposition together with the hypothesis $q_n^{r_n}\geqslant q_1^{n r_1}$, the same bound applies for the index of $V$. We conclude therefore that the index of $W$ is at most $\frac{1}{8} \epsilon^2 s$.

On the other hand, the index of the general element in the determinant for $W$ is at least
\[\phi_i - \sum_{m=1}^{n-1} j_m / r_m\]
where $\phi_i = \theta - i/r_n$, $\theta$ denotes the index of $P$, and
\[j_1+\dots+j_{n-1}=j\leqslant s-1\leqslant r_n;\]
further, by hypothesis, we have  $\delta r_m > r_{m+1}$ and so the above sum is at most $\delta$. Hence the index of $W$ is at least
\[\sum_{i=0}^{s-1} \max (\phi_i - \delta, 0) \geqslant \sum_{i=0}^{s-1} \max (\phi_i, 0) - s\delta.\]
But if $\theta r_n < s-1$ then the last sum is
\[([\theta r_n]+1)\left(\theta-[\theta r_n]/(2r_n)\right)\geqslant \tfrac{1}{4}\theta^2 s,\]
and if $\theta r_n \geqslant s-1$ then it is
\[\theta s - \tfrac{1}{2} s(s-1)/r_n \, \geqslant \, \tfrac{1}{2} \theta s.\]
On comparing estimates, we obtain
\[\max(\tfrac{1}{2}\theta, \tfrac{1}{4}\theta^2) \leqslant \frac{1}{8}\epsilon^2 + \delta \leqslant \frac{1}{4}\epsilon^2\]
whence $\theta \leqslant \epsilon$, as required.

\section{A combinatorial lemma}\label{sec:ch7_4}

We prove now a lemma of a combinatorial nature relating to the law of large numbers\index{law of large numbers}.\footnote{Cf. the paper of Wirsing\index{Wirsing, E.} cited earlier: Proc. Symposia Pure Math. (Amer. Math. Soc.), 20 (1971), 213--47.} A result of this kind occurred first in the work of Schneider\index{Schneider, Th.}, and it was utilized later by Roth\index{Roth, K. F.} who gave a simplified proof due to Davenport\index{Davenport, H.}. Another proof, attributed to Reuter\index{Reuter, G. E. H.}, and furnishing a slightly stronger theorem, was given by Mahler\index{Mahler, K.} in his tract, and Schmidt\index{Schmidt, W. M.} subsequently obtained the generalization we establish here.

\begin{mylemma}{}{ch7_2}
Suppose that  $r_1, ..., r_n$  and k are positive integers and that  $0 < \epsilon < 1$. Then the number of non-negative integers
\[j_{lm} \quad (1 \leqslant l \leqslant k, 1 \leqslant m \leqslant n)\]
satisfying
\[\sum_{l=1}^k j_{lm} = r_m \quad (1 \leqslant m \leqslant n), \qquad \sum_{m=1}^n j_{1m}/r_m < n/k - \epsilon n,\]
is at most
\[\binom{r_1+k-1}{r_1}\cdots\binom{r_n+k-1}{r_n}e^{-\frac{1}{4}\epsilon^2n}.\]
\end{mylemma}

\begin{proof}
We commence the proof by observing that the required number N of integers  $j_{lm}$  is given by  $\sum_{\nu_1(j_{11})\dots\nu_n(j_{1n})} \sum \nu_n(j_{2n})$ where the sum is over all non-negative integers  $j_1, ..., j_{1n}$  satisfying the given inequality, and  $\nu_m(j)$  denotes the number of solutions of the equation
\[\sum_{l=2}^{k} j_{lm} = r_m - j\]
in non-negative integers  $j_{2m},...,j_{km}$, that is
\[\nu_m(j) = \binom{r_m-j+k-2}{k-2}.\]
Hence we see that the multiple sum
\[\sum_{j_1=0}^{r_1} \dots \sum_{j_m=0}^{r_m} \nu_1(j_1) \dots \nu_m(j_m) \exp \left\{ \frac{1}{2} \varepsilon \left( n/k - \sum_{m=1}^n j_m/r_m \right) \right\}\]
is at least  $Ne^{\frac{1}{2}\epsilon^2n}$. Now the sum can be written alternatively in the form
\[\prod_{m=1}^{n} \left\{ \sum_{j_{m}=0}^{r_{m}} \nu_{m}(j_{m}) \exp\left(\frac{1}{2}\epsilon \rho_{m}\right) \right\},\]
where  $\rho_m = 1/k - j_m/r_m$, and clearly  $|\rho_m| \le 1$. But if  $|x| \le 1$  then  $e^x < 1 + x + x^2$, and so
\[\exp\left(\frac{1}{2}\epsilon\rho_m\right) < \frac{1}{2}\epsilon\rho_m + \exp\left(\frac{1}{4}\epsilon^2\right).\]
Further we have
\[\sum_{i=0}^{r_m} \nu_m(j_m) \rho_m = 0;\]
for  $\rho_m$  can plainly be expressed as
\[(r_m - j_m)/r_m - (1 - 1/k),\]
and it is easily verified by induction on r that
\[\sum_{j=0}^{r} \binom{r-j+k-2}{k-2} = \binom{r+k-1}{r},\]
\[\sum_{j=0}^r j \binom{j+k-2}{k-2} = r \left(1-\frac{1}{k}\right) \binom{r+k-1}{r}.\]
Thus, on appealing again to the first of the above binomial identities, we obtain
\[ \prod_{m=1}^{n} \left\{ \binom{r_m + k - 1}{r_m} e^{\frac{1}{4}\epsilon^2} \right\} \geqslant N e^{\frac{1}{4}\epsilon^2 n}, \]
and this gives the asserted estimate.
\end{proof}

\section{Grids}\label{sec:ch7_5}

Let T be a subspace of k-dimensional Euclidean space spanned by linearly independent vectors  $\mathbf{u}_1, ..., \mathbf{u}_{k-1}$. By a grid\index{grid} of size s on T we shall mean the finite set of vectors of the form
\[w_1\mathbf{u}_1+\ldots+w_{k-1}\mathbf{u}_{k-1},\]
where  $w_1, ..., w_{k-1}$  run through all rational integers with  $1 \le w_l \le s$.

Now let  $T_m$   $(1 \le m \le n)$  be any subspaces as above, and let  $\Gamma_m$  be a grid\index{grid} of size  $s_m$  on  $T_m$. Further let T,  $\Gamma$  signify the cartesian products  $T_1 \times \ldots \times T_n$  and  $\Gamma_1 \times \ldots \times \Gamma_n$  respectively. We shall denote by P a polynomial as indicated at the beginning of \S~3 with degree  $r_m$  in  $x_{1m}, \ldots, x_{km}$, and we shall signify by  $\Delta_m^{(j)}$  a differential operator as in \S~2, acting on  $x_{1m}, \ldots, x_{km}$. The following simple lemma, due to Schmidt\index{Schmidt, W. M.}, is fundamental to the proof of \autoref{thm:ch7_1}.

\begin{mylemma}{}{ch7_3}
If, for some integers  $t_m$   $(1 \le m \le n)$  with  $s_m(t_m+1) > r_m$, all polynomials  $\Delta_1^{(j_1)} \dots \Delta_n^{(j_n)} P$  with  $j_m \le t_m$  vanish everywhere on  $\Gamma$, then P vanishes identically on T.
\end{mylemma}

\begin{proof}
It is clear that the lemma follows at once by induction from the case n=1, and it will suffice therefore to prove the latter. Further, one can obviously assume, by applying a linear transformation, that  $T_m$  is the plane  $x_{km}=0$, with basis consisting of the first k-1 rows of the unit matrix of order k. Thus, omitting the suffix m, we see that it is enough to prove:

A polynomial  $P(x_1, ..., x_{k-1})$  with degree r vanishes identically if all  $\Delta^{(j)}P$  with  $j \leq t$  vanish at all integer points  $(w_1, ..., w_{k-1})$  with  $1 \leq w_l \leq s$, where s(t+1) > r.

Here  $\Delta^{(j)}$  denotes a differential operator on  $x_1, \ldots, x_{k-1}$  of order j. The assertion is clearly valid for k=2, since a polynomial in one variable with degree r cannot have more than r zeros, and we shall assume the proposition when k is replaced by k-1. If now P does not vanish identically then there is a largest integer q such that the rational function
\[ Q = (x_1 - 1)^{-q} \dots (x_1 - s)^{-q} P \]
is in fact a polynomial, and since, by hypothesis, s(t+1) > r, we have  $q \le t$. Further, by choice of q, one at least of the polynomials  $Q(w_1, x_2, ..., x_{k-1})$  with  $1 \le w_1 \le s$  does not vanish identically; let this be R. Then  $\Delta^{(j)}R$  vanishes at all integer points  $(w_2, ..., w_{k-1})$  with  $1 \le w_l \le s$, where  $\Delta^{(j)}$  is any differential operator on  $x_2, ..., x_{k-1}$ with order  $j \le t-q$. But R has degree at most r-sq < (t-q+1)s, and this is plainly contrary to the inductive hypothesis. The contradiction establishes the assertion.
\end{proof}

\section{The auxiliary polynomial}\label{sec:ch7_6}

For each m with  $1 \le m \le n$  we shall denote by  $L_{lm}$  ( $1 \le l \le k$ ) linear forms in  $x_{1m}, \ldots, x_{km}$  with real algebraic integer\index{algebraic integer} coefficients. Further we shall denote by d the degree of the field K generated by all the coefficients over the rationals, and we shall signify by  $c_1, c_2, \ldots$  numbers\index{numbers} greater than 1 which depend on these coefficients only.

Let now  $r_1, ..., r_n$  be any positive integers, and let  $r = \max r_m$. Further suppose that  $0 < \epsilon < 1$  and that  $e^{\frac{1}{4}\epsilon^2n} > 2kd$. Adopting the notation of \autoref{sec:ch7_3}, we have

\begin{mylemma}{}{ch7_4}
There is a polynomial P with degree at most  $r_m$  in  $x_{1m}, \ldots, x_{km}$  and with height at most  $c_1^r$  such that, for each l with  $1 \le l \le k$, the index of P with respect to the  $L_{lm}$  and  $r_m$  is at least  $n/k - \varepsilon n$.
\end{mylemma}

\begin{proof}
It can be assumed, without loss of generality, that, for all l, m, the coefficient of  $x_{1m}$  in  $L_{lm}$, say  $\alpha_{lm}$, is not 0. Then P has to be determined such that, for all l and all non-negative integers  $j_1, \ldots, j_n$  with
\[\sum_{m=1}^{n} j_m/r_m < n/k - \epsilon n,\]
the polynomials
\[(j_1! \dots j_n!)^{-1} (\partial/\partial x_{11})^{j_1} \dots (\partial/\partial x_{1n})^{j_n} P\]
vanish identically when  $-L_{lm}$, with  $x_{1m}$  equated to 0, is substituted for  $x_{1m}$, and the factor  $\alpha_{lm}$  is included to multiply each of  $x_{2m}, \ldots, x_{km}$. Now these polynomials are homogeneous in  $x_{2m}, \ldots, x_{km}$  with degree  $r_m - j_m$  and hence, by \autoref{lem:ch7_2}, they have, in total, at most  $kNe^{-\frac{1}{4}e^2n}$  coefficients, where N denotes the product of binomial factors occurring in the enunciation of the lemma. Each coefficient is a linear form in the coefficients of P, and there are precisely N of the latter. Furthermore, the coefficients in the linear forms are algebraic integers\index{algebraic integer} in K with sizes at most  $c_2^r$  (cf. the estimates in \autoref{sec:ch7_3}). It follows, on utilizing an integral basis for K and recalling the hypothesis  $e^{\frac{1}{4}e^2n} > 2kd$, that one has to satisfy at most  $\frac{1}{2}N$  linear equations\index{linear equations} with rational integer coefficients each having absolute value at most  $c_3^r$  (cf. \S~3, \autoref{ch:6}). The required result is now obtained from \autoref{lem:ch6_1} of \autoref{ch:6}.
\end{proof}

\section{Successive minima}\label{sec:ch7_7}

We recall from the Geometry of Numbers\index{Geometry of Numbers}\index{numbers} that if R is any convex body in k-dimensional Euclidean space, then the numbers\index{numbers}  $\lambda_l$   $(1 \le l \le k)$, given by the infimum of all  $\lambda > 0$  such that  $\lambda R$  contains l linearly independent integer points, are termed the successive minima\index{successive minima} of R, and they have the property that  $\lambda_1 \dots \lambda_k V$, where V denotes the volume of R, is bounded above and below by positive numbers\index{numbers} depending only on k.

We now combine the preceding lemmas to obtain a proposition on the penultimate minimum of a certain parallelepiped, which will be the main instrument in the proof of \autoref{thm:ch7_1}. We shall denote by  $M_1, \ldots, M_k$  linear forms in  $x_1, \ldots, x_k$  with real algebraic integer\index{algebraic integer} coefficients constituting a non-singular matrix A, and we shall signify by  $M'_1, \ldots, M'_k$  the adjoint linear forms with coefficients given by the columns of  $A^{-1}$. Further we shall signify by S some non-empty set of suffixes i such that  $M'_i$  does not represent zero for any integral values, not all 0, of the variables; the assumption that S exists involves, of course, some loss of generality. We prove:

\begin{mylemma}{}{ch7_5}
For any  $\zeta > 0$  there exists c > 0 such that for all positive  $\mu_1, \ldots, \mu_k$  satisfying  $\mu_1 \ldots \mu_k = 1$  and  $\mu_i \ge 1$  for i in S, the penultimate minimum  $\lambda_{k-1}$  of the parallelepiped  $|M_i| \le \mu_i$   $(1 \le l \le k)$  exceeds  $\mu^{-l}$, where  $\mu$  denotes the maximum of  $\mu_1, \ldots, \mu_k$  and c.
\end{mylemma}

\begin{proof}
It will be seen that the lemma immediately implies Roth\index{Roth, K. F.}'s theorem\index{Roth's theorem}, that is the case n = 1 of \autoref{thm:ch7_1}; this follows on taking
\[M_1 = \alpha_1 x_1 - x_2, \quad M_2 = x_2\]
and S to consist just of the suffix 2, as is possible since  $\alpha_1$  is irrational.

We show first that it suffices to prove a modified version of \autoref{lem:ch7_5}. Suppose that  $Q \geqslant \mu^k$  and let  $\omega_1, \ldots, \omega_k$  be defined by  $\mu_l = Q^{\omega_l}$. Then since  $\mu_1 \ldots \mu_k = 1$  we have  $\omega_1 + \ldots + \omega_k = 0$  and clearly  $\omega_i \geqslant 0$  for i in S. Clearly also  $\omega_l \leqslant 1$  for all l and, since again  $\mu_1 \ldots \mu_k = 1$  we have  $\mu_l \geqslant Q^{-1}$, whence  $\omega_l \geqslant -1$. Now for any positive integer N there are rationals  $\omega_1', \ldots, \omega_n'$  with denominator N satisfying  $|\omega_l - \omega_l'| < 1/N$  and  $|\omega_l'| \leqslant 1$  for all l, and also  $\omega_1' + \ldots + \omega_k' = 0$; indeed one has merely to take  $N\omega_1' = [N\omega_1]$  and, having defined  $\omega_1', \ldots, \omega_{l-1}'$, to take  $N\omega_l'$  as  $[N\omega_l]$  or  $-[-N\omega_l]$  according as  $\omega_1' + \ldots + \omega_{l-1}'$  does or does not exceed  $\omega_1 + \ldots + \omega_{l-1}$. Plainly the  $\omega_1', \ldots, \omega_k'$  belong to a finite set of rationals independent of Q, and the minimum  $\lambda_{k-1}'$  of the parallelepiped  $|M_l| \leqslant Q^{\omega_l'}$  ( $1 \leqslant l \leqslant k$ ) exceeds  $Q^{-1/N}\lambda_{k-1}$. Hence it is enough to prove:

For any real  $\omega_1, \dots, \omega_k$  with  $\omega_1 + \dots + \omega_k = 0$,  $|\omega_l| \le 1$   $(1 \le l \le k)$  and  $\omega_i \ge 0$  for all i in S, and for any  $\zeta > 0$, there exists C > 0 such that, for all Q > C, the minimum  $\lambda_{k-1}$  of the parallelepiped
\[|M_l| \leqslant Q^{\omega_l} \quad (1 \leqslant l \leqslant k)\]
exceeds  $Q^{-\zeta}$.

We shall suppose that  $\zeta \leq 1$, as obviously we may, and we shall signify by d the degree of the field generated by the elements of A over the rationals. Let  $\varepsilon$  be any positive number less than  $\zeta/(8k)^2$, let n be any integer satisfying the condition preceding \autoref{lem:ch7_4}, and let  $\delta$  be defined as in \autoref{lem:ch7_1}. We shall assume that there is an unbounded set of values of Q such that  $\lambda_{k-1} \leq Q^{-\zeta}$, and we shall ultimately derive a contradiction. We select a sufficiently large  $Q_1$  in this set, that is  $Q_1 > c_1$, where  $c_1$, like  $c_2$,  $c_3$  below, depends only on A, k, n, d,  $\varepsilon$,  $\delta$,  $\zeta$  and the  $\omega$'s. We then select further elements  $Q_2, \ldots, Q_n$  in the set such that  $Q_m^{\frac{1}{2}\delta} > Q_{m-1}$   $(1 < m \leq n)$, whence clearly  $Q_1 < \dots < Q_n$. Finally we choose positive integers  $r_1, \ldots, r_n$  such that  $Q_1^{\text{tr}_1} > Q_n$  and
\[Q_1^{r_1}\leqslant Q_m^{r_m}\leqslant Q_1^{(1+\epsilon)r_1}\quad (1\leqslant m\leqslant n);\]
then plainly the condition preceding \autoref{lem:ch7_1} is satisfied.

We observe now that the hypotheses of \autoref{lem:ch7_4} hold when $Ifa =$ $M_{lm}(x_m)$, where $x_m$ denotes the vector $(x_{1m},\dots,x_{km})$; let P be the polynomial constructed there. Further we note that, for any Q as above, there exist linearly independent integer points $\mathbf{u}_1, \dots, \mathbf{u}_k$ with $\mathbf{u}_l$ in $\lambda_l R$, where R denotes the given parallelepiped and $\lambda_1, \dots, \lambda_k$ its successive minima\index{successive minima}. Moreover, there is a linear form L with relatively prime integer coefficients, unique except for a factor + 1, which vanishes at $\mathbf{u}_1, \dots, \mathbf{u}_{k-1}$; we take $\mathbf{u}_{lm}$ and $L_m$ to be these $\mathbf{u}_l$ and L respectively when $Q = Q_m$. We shall verify later that, if Q is sufficiently large, then q = ||L|| satisfies $Q^c < q < Q^{c'}$, where c, c' are positive numbers\index{numbers} depending only on $\zeta$ and d. Assuming this for the present, it follows that all the hypotheses of \autoref{lem:ch7_1} are satisfied with $v = c/c'$, provided that $c_1$, and so also $q^1$ and $Q_1$ are large enough. Hence we conclude that the index of P with respect to the $L_m$ and $r_m$ is at most $\varepsilon$.

We proceed to prove that, with the notation of \autoref{sec:ch7_5}, all polynomials $\Delta P$ with
\[\Delta = \Delta_1^{(j_1)} \dots \Delta_n^{(j_n)}, \quad \sum_{m=1}^n j_m/r_m < 2\epsilon n\]
vanish everywhere on $\Gamma$, where $\Gamma_m$ is the grid\index{grid} of size $[\epsilon^{-1}] + 1$ on the space $T_m$ spanned by $\mathbf{u}_{lm}$ ($1 < l < k$). This implies, by \autoref{lem:ch7_3}, on taking $t_m = [\epsilon r_m]$, that all polynomials $\Delta P$, with $\sum j_m/r_m < \epsilon n$, vanish identically on the $n(k-1)$-dimensional space of solutions of
\[L_1 = \dots = L_n = 0.\]
But the latter contradicts the above conclusion concerning the index of P, and the contradiction establishes the lemma. To prove the proposition, let $\Delta P$ be any of the polynomials in question and let P' be the polynomial in new variables $y_{lm}$ obtained from $\Delta P$ by the linear substitution $y_{lm} = X_{jm}$. Then it is readily verified that P' has height at most $c_1^r$, where $r = \max r_m$. Further, since, by assumption, $\lambda_{k-1} < Q^{-\zeta}$, we have for any $x_m$ on $\Gamma_m$
\[|y_{lm}| < k(\epsilon^{-1}+1) Q_m^{\omega_l-\zeta} < Q_m^{\omega_l-\frac{1}{2}\zeta}.\]
Thus, by \autoref{lem:ch7_4}, it follows that, for all points on $\Gamma$, we have $|\Delta P| < c_2^r e^s$, where
\[s = \sum_{l=1}^{k} \sum_{m=1}^{n} (\omega_{l} - \frac{1}{2}\zeta) j_{lm} \log Q_{m},\]
and $j_{lm}$ are some non-negative integers with
\[\sum_{l=1}^{k} j_{lm} \leqslant r_{m} \quad (1 \leqslant m \leqslant n), \quad n/k - \sum_{m=1}^{n} j_{lm}/r_{m} < 3\epsilon n \quad (1 \leqslant l \leqslant k).\]
Denoting, for brevity, the left-hand side of the last inequality by  $h_l$, we see that, by the first inequality,  $h_1 + \ldots + h_k \ge 0$, and so both inequalities together imply that  $|h_l| < 3kne\epsilon$. Further, since  $|\omega_l| \le 1$, we obtain, in view of the initial choice of  $r_1, \ldots, r_n$,
\[s\leqslant r_1\log Q_1\sum_{l=1}^k\sum_{m=1}^n\{(\omega_l-\tfrac12\zeta)j_{lm}/r_m+2\epsilon\}.\]
But now, by virtue of our estimate for  $h_l$  and the hypothesis
\[\omega_1 + \ldots + \omega_k = 0,\]
the double sum here differs from  $-\frac{1}{2}\zeta n$  by at most  $8k^2n\epsilon$. Since, by definition,  $\epsilon < \zeta/(8k)^2$, it follows that  $|\Delta P| < Q_1^{-\frac{1}{4}\zeta nr} < 1$, provided  $Q_1$  is sufficiently large. On the other hand,  $\Delta P$  is a rational integer for all points on  $\Gamma$, and hence  $\Delta P = 0$, as required.

It remains only to prove the assertion concerning q = ||L||. Let U be the matrix with columns  $\mathbf{u}_1, ..., \mathbf{u}_k$  and let  $\mathbf{v}_1, ..., \mathbf{v}_k$  be the rows of  $U^{-1}$. Then clearly  $\rho \mathbf{v}_k$  is the coefficient vector of L for some rational  $\rho$. Since  $L(\mathbf{u}_k)$  is an integer and  $\mathbf{v}_k \mathbf{u}_k = 1$,  $\rho$  is in fact an integer. Further  $\rho$ divides det U, for plainly  $U^{-1} = \operatorname{adj} U/\operatorname{det} U$. Furthermore we have det  $U \leq 1$, where the implied constant depends only on A, for certainly R has volume  $\gg 1$  and hence, by the property of successive minima\index{successive minima} quoted at the beginning,  $\det(AU) \leqslant 1$. It follows that each element of  $(\det \mathbf{U}) \mathbf{v}_k$  is a rational integer  $\leqslant q$. Hence the element in the kth row and lth column of adj (AU), namely  $(\det(AU)) M'_{l}(\mathbf{v}_{k})$, is an algebraic integer\index{algebraic integer} with size  $\leq q$. But by hypothesis we have  $\lambda_{k-1} < Q^{-\zeta}$  and  $\omega_1 + \ldots + \omega_k = 0$, and thus the element is also  $\ll Q^{-\omega_l-(k-1)\zeta}$. We conclude that, for l in S, the element is both  $\Rightarrow q^{-d}$\footnote{We are using Vinogradov\index{Vinogradov, A. I.}’s notation; by $a \ll b$ or $a \leqslant b$ one means $|a| < cb$ for some constant c, and similarly for $\gg$.} and  $\leqslant Q^{-(k-1)\zeta}$, and, since S is assumed non-empty, this gives the required lower bound for q. The upper bound follows from the identity  $U^{-1} = (AU)^{-1} A$, on observing, as above, that the elements in the kth row of  $(AU)^{-1}$  are  $\leqslant Q$.
\end{proof}

\section{Comparison of minima}\label{sec:ch7_8}

We prove first a general lemma of Davenport\index{Davenport, H.}, and we proceed then to show that, with some proviso, the minima  $\lambda_{k-1}$  and  $\lambda_k$  of the parallelepiped of \autoref{lem:ch7_5} differ only by a small factor. Constants implied by  $\leqslant$  will depend only on k.

\begin{mylemma}{}{ch7_6}
Let  $L_1, ..., L_k$  be real linear forms with determinant 1 and let  $\lambda_1, ..., \lambda_k$  be the successive minima\index{successive minima} of the parallelepiped
\[|L_l| \leqslant 1 \quad (1 \leqslant l \leqslant k).\]
Suppose that  $\rho_1 \ge ... \ge \rho_k > 0$  and that
\[\rho_1 \lambda_1 \leqslant \ldots \leqslant \rho_k \lambda_k, \quad \rho_1 \ldots \rho_k = 1.\]
Then for some permutation  $\rho'_1, \ldots, \rho'_k$  of  $\rho_1, \ldots, \rho_k$, the successive minima\index{successive minima}  $\lambda'_1, \ldots, \lambda'_k$  of the parallelepiped  $\rho'_l | L_l | \leq 1$   $(1 \leq l \leq k)$  satisfy
\[\rho'_l \lambda_l \leqslant \lambda'_l \leqslant \rho'_l \lambda_l \quad (1 \leqslant l \leqslant k).\]
\end{mylemma}

\begin{proof}
There certainly exist linearly independent integer points  $\mathbf{x}_1, \ldots, \mathbf{x}_k$  such that one at least of  $|L_1|, \ldots, |L_k|$  assumes the value  $\lambda_l$  at  $\mathbf{x}_l$, and we denote by  $S_l$  the space spanned by  $\mathbf{x}_1, \ldots, \mathbf{x}_l$. Further, for each  $l \geq 2$, there is a non-trivial linear relation  $a_1 L_1 + \ldots + a_l L_l = 0$  satisfied identically on  $S_{l-1}$, and  $L_1, \dots, L_k$  can be permuted so that  $|a_l|$  is maximal; this gives
\[|L_1| + \ldots + |L_{l-1}| > \frac{1}{2}(|L_1| + \ldots + |L_l|)\]
identically on $S_l$ for $l = 1, 2, \dots, k$. Now for any j it is clear that every point in $S_l$ not in $S_{j-1}$ satisfies
\[\max (|L_1|, \dots, |L_l|) > \lambda_l,\]
and thus, in view of the inequality obtained above, it satisfies also
\[\max (\rho_1 |L_1|, \dots, \rho_k |L_k|) > \rho_j \lambda_l.\]
By hypothesis, $\rho_j \lambda_j \le \rho_l \lambda_l$ for $j > l$, and the required lower bound for $\lambda'_l$ follows on taking $\rho'_1, \dots, \rho'_k$ to be the permutation of $\rho_1, \dots, \rho_k$ inverse to that associated with $L_1, \dots, L_k$. The upper bound is a consequence of the equation $\rho'_1 \dots \rho'_k = 1$ together with the property, noted earlier, that $\lambda_1 \dots \lambda_k$ and $\lambda'_1 \dots \lambda'_k$ are both $\ll 1$ and $\gg 1$.
\end{proof}

\begin{mylemma}{}{ch7_7}
The last two minima of the parallelepiped of \autoref{lem:ch7_5} satisfy $\lambda_{k-1} > \lambda_k \mu^{-k^2\epsilon}$, provided that $\lambda_1 \mu^4 > \mu^{-\epsilon}$ for all i in S.
\end{mylemma}

\begin{proof}
The hypotheses of \autoref{lem:ch7_6} hold with $L_l$ ($1 \le l \le k$) given by $M_l/\mu_l$ and $\rho_l$ given by $\mu_l \mu^{-1/k}$ ($1 \le l \le k$), where $\mu$ is defined by the equation $\mu_1 \dots \mu_k = 1$. Let $\rho'_1, \dots, \rho'_k$ be the permutation of $\rho_1, \dots, \rho_k$ indicated in the lemma, and let $\mu'_l = \mu_l \mu^{-1/k}/\rho'_l$. Assume first that $\mu'_i > 1$ for all i in S. Then from \autoref{lem:ch7_5} with $\mu_l$ replaced by $\mu'_l$, we infer that, for any $\epsilon' > 0$, there exists $c' > 0$ such that $\lambda'_{k-1} > \mu'^{-\epsilon'}$, where $\mu'$ denotes the maximum of $\mu'_1, \dots, \mu'_k$ and $c'$. On the other hand, from \autoref{lem:ch7_6}, $\lambda'_{k-1} < \rho'_{k-1} \lambda_{k-1}$ and clearly, since $\lambda_1 \dots \lambda_k \ll 1$, we have $\rho'_{k-1} \ll \lambda_{k-1}/\lambda_k$. Thus it suffices to prove that $\mu'^{\epsilon'} \ll \mu^\epsilon$ if $\epsilon'$ is chosen sufficiently small. But by hypothesis, since S is assumed non-empty, we have $\lambda_1 \mu > \mu^{-\epsilon}$; further, since $\lambda_l > \lambda_1$ for all l, we see that $\mu > \lambda_1$ and $\mu_l^{-1} \lambda_{k-1} < 1$. Hence we obtain
\[\mu'_l < \lambda_{k-1}\mu < \mu^{k^2\epsilon+1},\]
and the required result follows. If, contrary to the above assumption, $\mu'_i < 1$ for some i in S, then, on observing that by hypothesis we obtain $\mu > \mu^{-\epsilon}$ and the required result again follows.
\end{proof}

\section{Exterior algebra}\label{sec:ch7_9}

For any vectors $\mathbf{x}_1, \dots, \mathbf{x}_l$ in $R^k$ with $1 < l < k$, one denotes by $\mathbf{x}_1 \wedge \dots \wedge \mathbf{x}_l$ the vector in $R^m$ whose elements are the $m = \binom{k}{l}$ subdeterminants of order l formed from the k by l matrix with columns $\mathbf{x}_1, \dots, \mathbf{x}_l$. We shall utilize some well-known properties of this product; in particular, we shall require Laplace's identity\index{Laplace's identity}
\[( \mathbf{x}_1 \wedge \dots \wedge \mathbf{x}_l ) \cdot ( \mathbf{y}_1 \wedge \dots \wedge \mathbf{y}_l ) = \operatorname{det}(\mathbf{x}_i \cdot \mathbf{y}_j),\]
where on the left one has the usual vector dot product, and also the relation
\[\operatorname{det} M_\sigma = (\operatorname{det} A)^{m-1},\]
which holds for any matrix A of order k with column vectors $\mathbf{a}_1, \dots, \mathbf{a}_k$, say, where $M_\sigma = \mathbf{a}_{i_1} \wedge \dots \wedge \mathbf{a}_{i_l}$, and $\sigma$ runs through all sets of l distinct integers $i_1, \dots, i_l$ from $1, \dots, k$.\footnote{Short proofs are given in Schmidt\index{Schmidt, W. M.}'s tract (Bibliography).}

We shall need, in addition, the following lemma, due to Mahler\index{Mahler, K.}, on compound convex bodies. A will signify a matrix as above with det A = 1, and, as in \autoref{sec:ch7_8}, constants implied by $\ll$ will depend only on A. Further we shall denote by $\mathbf{a}\mathbf{x}$ the linear form in the elements of $\mathbf{x}$ with coefficient vector $\mathbf{a}$.

\begin{mylemma}{}{ch7_8}
The successive minima\index{successive minima} $\lambda_1, \dots, \lambda_k$ and $\nu_1, \dots, \nu_m$ of the parallelepipeds $|\mathbf{a}_i \mathbf{x}| \le 1$ ($1 \le i \le k$) and $|M_\sigma \mathbf{X}| \le 1$ respectively, where $\sigma$ runs through all sets $\sigma$ as above, satisfy $\nu_\sigma \ll \lambda_\sigma$, where $\lambda_\sigma = \prod \lambda_j$, the product being taken over all j in $\sigma$, and $\nu_1 \le \nu_2 \le \dots \le \nu_m$.
\end{mylemma}

\begin{proof}
Let $\mathbf{x}_1, \dots, \mathbf{x}_k$ be linearly independent integer points such that $|\mathbf{a}_i \mathbf{x}_j| \le \lambda_j$ ($1 \le j \le k$), and let $\mathbf{X}_\sigma$ be defined like $\lambda_\sigma$ above, with $\mathbf{x}$ in place of $\mathbf{a}$. By Laplace's identity\index{Laplace's identity} we have
\[|M_\sigma \mathbf{X}_\tau| = |\operatorname{det}(\mathbf{a}_i \mathbf{x}_j)| \leqslant m \lambda_\sigma,\]
where i, j run through all elements of $\sigma$, $\tau$ respectively. Hence, for each i with $1 \leqslant i \leqslant m$, we have $|M_\sigma \mathbf{X}_i| \ll \lambda_\sigma$ and so $\nu_i \ll \lambda_\sigma$. But since, by hypothesis, det A = 1, we have det $M_\sigma = 1$, and thus the volume of the parallelepiped $|M_\sigma \mathbf{X}| \leqslant 1$ is $2^m$. Thus $\nu_1 \dots \nu_m \gg 1$ and since $\lambda_{\sigma_1} \dots \lambda_{\sigma_m} \ll 1$ it follows that $\nu_i \gg \lambda_\sigma$, as required.
\end{proof}

\section{Proof of main theorem}\label{sec:ch7_10}

It will suffice to prove \autoref{thm:ch7_1} under the assumption that $\alpha_1, \dots, \alpha_n$ are real algebraic integers\index{algebraic integer}, for clearly the general result then follows on multiplying each $\alpha_i$ by the leading coefficient in its minimal polynomial\index{minimal polynomial of algebraic number}. We shall signify by $\mathbf{a}_j$ ($1 \leqslant j \leqslant n$) the vector in $R^{n+1}$ given by $(\mathbf{e}_j, \alpha_j)$, where $\mathbf{e}_1, \dots, \mathbf{e}_n$ denote the rows of the unit matrix of order n. Further, for brevity, we shall write k = n + 1, and we shall denote by $\mathbf{a}_k$ the vector (0, ..., 0, 1) in $R^k$. Constants implied by $\ll$ or $\gg$ will depend only on $\alpha_1, \dots, \alpha_n$ and the quantities $\epsilon, \zeta$ to be defined below.

We show first that the theorem is a consequence of the following proposition.

For any $\zeta > 0$ and any positive numbers\index{numbers} $\mu_1, \dots, \mu_k$ with $\mu_1 \dots \mu_k = 1$ and $\mu_j \le 1$ ($1 < j < k$), the first minimum $\lambda_1$ of the parallelepiped $|\mathbf{a}_j \mathbf{x}| \le \mu_j$ ($1 < j < k$) exceeds $\mu^{-1}$ if $\mu > \mu_k$ and $\mu \gg 1$.

The proof proceeds by induction on n; we have already remarked that the case n = 1 is an immediate consequence of \autoref{lem:ch7_5}, and we assume now that the theorem holds when n is replaced by n — 1. Let q be a positive integer satisfying the inequality occurring in the enunciation, and let $\mu_j = q^{-1-1/n-\epsilon}$ ($1 \le j \le n$). Further let $\mu_k = (\mu_1 \dots \mu_n)^{-1}$, where k = n + 1 as above. Then clearly $\mu_k > q^{1+\epsilon'}$ and moreover the first minimum $\lambda_1$ of the parallelepiped $|\mathbf{a}_j \mathbf{x}| < \mu_j$ ($1 < j < k$) is at most $q^{-\epsilon'/n}$. But, on appealing again to the given inequality and applying the inductive hypothesis, we see that, if $q \gg 1$, then $\mu_j < 1$ for all $j < k$. Hence the proposition above shows that $\lambda_1 > \mu^{-\epsilon}$ for any $\epsilon > 0$ and any $\mu$ with $\mu \gg 1$ and $\mu \ge \mu_k$. Furthermore, by the case n = 1 of the theorem, we have $\mu_j > q^{-1}$ ($1 \leqslant j < n$), whence $\mu_k < q^n$. Plainly the estimates for $\lambda_1$ are inconsistent if $\epsilon$ is sufficiently small, and the contradiction proves the theorem.

Preliminary to the proof of the proposition, 
we observe that, with the notation of \autoref{sec:ch7_9}, the linear forms $M_\sigma = M_\sigma \mathbf{X}$ satisfy the hypotheses of \autoref{sec:ch7_7} with S given by those sets $\sigma$ which include k. For it is easily verified from the Laplace expansions of A that, as $\sigma$ runs through the complement of $\sigma$ in $1, \dots, k$, the forms $M'_\sigma \mathbf{X}$ constitute the set adjoint to the $M_\sigma$, except possibly for a sign change; further, if $\sigma$ does not include k, we have
\[M'_\sigma \mathbf{X} = \mathbf{X}_\sigma + \sum ( \pm \alpha_j ) \mathbf{X}_{\sigma-j+k},\]
where the summation is over all j in $\sigma$, on the right there occur the co-ordinates of \textbf{X}, and $\sigma-j+k$ denotes the set $\sigma$ with k in place of j. By hypothesis $1, \alpha_1, \dots, \alpha_n$ are linearly independent over the rationals, and thus we see that $M'_\sigma \mathbf{X} \neq 0$ for all integer vectors $\mathbf{X} \neq 0$, as required.

The proof of the proposition proceeds by induction on k; 
the result plainly holds for $k = 2$ by \autoref{lem:ch7_5}, 
and we assume now that it has been verified for all values up to $k — 1$. 
Let $l$ be any integer with $1 \le l < k$ and, 
for any set $\tau$ of $l$ distinct integers from $1, \dots, k$, 
let $\mu_\tau = \prod \mu_j$, where the product is over all j in $\tau$. 
By \autoref{lem:ch7_7} we see that the successive minima\index{successive minima} $\nu_1, \dots, \nu_m$ of the 
parallelepiped $|M_\tau| < \mu_\tau$ satisfy $\nu_{m-1} \gg \nu_m \mu^{-k\zeta}$ for any $\tau>0$, 
provided that $\mu \gg 1$, $\mu \gg \mu_\tau$ for all $\tau$ 
and $\nu_1 \mu_\tau > \mu^{-\zeta}$ for $\tau$ in S. 
Further, with the notation of \autoref{lem:ch7_8}, 
it is clear that $\tau_{m}$ and $\tau_{m-1}$ consist of the integers $k — l+1, k - l + 2, ..., k$ and $k — l, k — l + 2, ..., k$ respectively. 
Thus, under the above conditions, 
we have 
\[
    \lambda_{k-1}\lambda_{k-l+2}\cdots\lambda_k \gg \lambda_{k-l+1}\cdots\lambda_k \mu^{-k^2\epsilon}\mu^{-k\zeta} 
\] 

that is $\lambda_{k-l} \gg \lambda_{k-l+1} \mu^{-k\zeta}$. 
The required inequality $\lambda_1 > \mu^{-\zeta}$ follows on applying the latter with $l = 1, 2, ..., k — 1$, 
noting that $\lambda_k \gg 1$, and taking $\zeta$ sufficiently small.

Since evidently $\mu_\tau \le \mu_k$ for all $\tau$, 
it remains only to prove that, for $\tau$ in S, 
$\nu_1 \mu_\tau > \mu^{-\zeta}$ for any $\mu$ with $\mu \gg 1$ 
and $\mu > \mu_k$. In fact it suffices to show that $\lambda_\sigma \mu_\sigma > \mu^{-\zeta}$, 
for, again from \autoref{lem:ch7_8}, we have $\nu_1 \gg \lambda_1 \dots \lambda_l \geqslant \lambda_1^{l}$. 
Now, by the definition of $\lambda_l$, 
the parallelepiped $|\mathbf{a}_j \mathbf{x}| < \lambda_l/\mu_j$ ($1 \leqslant j \leqslant k$) contains an 
integer point $\mathbf{x} \neq 0$; 
in fact the $k$th co-ordinate of \textbf{x} is not 0 since $\lambda_1 \leqslant 1$ 
by Minkowski's linear forms theorem\index{Minkowski's linear forms theorem}, 
whence 
$\lambda_1\mu_j <  1(1\leqslant j \leqslant k)$, 

and $\mathbf{a}_j \mathbf{x}$ is simply the jth co-ordinate of \textbf{x} when the $k$th co-ordinate vanishes. 
It follows that, if $\tau$ is any element of S, 
then the parallelepiped in $R^l$ given by $|\mathbf{a}_i \mathbf{x}| \leqslant \lambda_l/\mu_i$, 
where $i$ is restricted to $\tau$ and the co-ordinates of \textbf{x} with suffixes not in $\tau$ 
are disregarded, also contains an integer point $\mathbf{x} \neq 0$. 
Hence the first minimum $\lambda_1'$ of the parallelepiped $|\mathbf{a}_i \mathbf{x}| \le \mu_i'$ in $R^l$, 
where $\mu_i' = \mu_i/\lambda_l$, 
is at most $\lambda_l/\mu_\tau^{1/l}$. 
It is therefore enough to prove that $\lambda_1' > \mu^{-\zeta}$; 
but this follows from the inductive hypothesis since clearly $\prod \mu_i' = 1$ and $\mu'_\sigma > 1$. 
The theorem is herewith established.
